\input{@resources/tex-add-search-paths}  % allow latex to find custom stuff
\input{./.econtexRoot}                   % Set path to root dir
\documentclass{econark}
\usepackage{econark}                     % shared notation and shortcuts
\usepackage{econark-notation}            % the registry's own macros
\usepackage{llorracc-handouts}           % \handoutHeader, \handoutNameMake
\providecommand{\texname}{RamseyCassKoopmans}% Keyname for the paper
\providecommand{\FigDir}{}\renewcommand{\FigDir}{\econtexRoot/Figures}
% The notebook that solves the model and draws both figures. The handout keeps its own copy beside it,
% identical to the book's beside the chapter (Chris, 2026-11-14, in cc:course-choice: "own copy of the
% notebook"), and links it as a file in the same folder, as the book's [notebook](RamseyCassKoopmans.ipynb) does.
\providecommand{\Notebook}{RamseyCassKoopmans.ipynb}
\providecommand{\NotebookColab}{https://colab.research.google.com/github/llorracc/RamseyCassKoopmans/blob/master/RamseyCassKoopmans.ipynb}
\providecommand{\NotebookBinder}{https://mybinder.org/v2/gh/llorracc/RamseyCassKoopmans/master?labpath=RamseyCassKoopmans.ipynb}
% Only the book hosts JupyterLite, so this link still runs the book's copy.
\providecommand{\BookNotebookLite}{https://intertemporal-choice.github.io/lite/notebooks/index.html?path=RamseyCassKoopmans.ipynb}
                 % macros for various subdirs
\usepackage{@local/\texname}             % macros specific to this handout
% endfloat defines \filename, which \bibfilesfound reads when \bibliography expands it; so load it first
\usepackage[noheads,nolists,nomarkers]{endfloat}
\renewcommand{\efloatheading}[1]{}       % no "FIGURES" running head on the figure pages, as in the 2020 handout
\usepackage{econark-bibfilesfind}        % Package to find valid bibfiles
\bibfilesfind{\texname}                  % find them
\usepackage{hyperref}                    % allow hyperlinks
\usepackage{econark-make4ht}             % link colors and pdf-or-web settings

% The book chapter is the source of this text (content/growth/RamseyCassKoopmans/index.md in
% intertemporal-choice.github.io). Make edits there first and carry them here; check-book-sync.py
% reports any differences. The numerical appendix is split off: the time-elimination method now lives in
% the notebook (RamseyCassKoopmans.ipynb, the book's copy), as in the chapter.
% A good follow-up handout (or exam question) would look at the augmented model in \cite{ghosDespiteUzawa}
\begin{document}
\handoutHeader

\begin{verbatimwrite}{\jobname.title}
  The Ramsey/Cass-Koopmans (RCK) Model
\end{verbatimwrite}

\handoutNameMake


\citet{ramseySave}, followed much later by \citet{cass:growth} and \citet{koopmans:growth}, formulated the canonical model of optimal growth, presented here for an economy with exogenous ``labor-augmenting'' technological progress. For the facts the model is meant to explain, and the Solow model it builds on, see the undergraduate textbook of \citet{ashrafWeilGrowth}.\marginFootnote{Graduate textbook treatments are in \citet{blanchard&fischer:text}, chapter 2, \citet{barro&salaimartin:book}, chapter 2, and \citet{acemogluGrowth}, chapter 8.}

\section{The Budget Constraint}

The economy has a perfectly competitive production sector that uses a constant-returns-to-scale Cobb-Douglas production function\marginFootnote{States and controls are functions of time; time subscripts appear only where clarity needs them.}
\begin{equation}\begin{gathered}\begin{aligned}
                   \YLvl & =  \FFunc(\KLvl,\PtyLab \LLvl) = \KLvl^{\kapShare}(\PtyLab \LLvl)^{1-\kapShare}
\end{aligned}\end{gathered}\end{equation}
to produce output using capital and labor.  Under perfect competition firms must make zero profits in equilibrium, which means, by \href{https://intertemporal-choice.github.io/content/appendices/mathematical/mathfactslist/\#fact-eulerstheorem}{Euler's theorem}, that\marginFootnote{An important insight here is that all of output is absorbed by payments to labor and to capital; there is nothing left over (say, to fund corporate research to invent new technologies to augment productivity). \href{https://intertemporal-choice.github.io/content/growth/endogenous/}{What Endogenous Growth Requires} starts from this point.}
\begin{equation}
  \YLvl = \FFunc_{\KLvl}\KLvl + \FFunc_{\LLvl}\LLvl , \qquad \FFunc_{\KLvl} \equiv \partial\YLvl/\partial\KLvl , \quad \FFunc_{\LLvl} \equiv \partial\YLvl/\partial\LLvl . \label{eq:RCK-outputExhausted}
\end{equation}
Labor hours (the same as population) increase exogenously at a constant rate\marginFootnote{The population growth rate is written $\PopnGroRte$: $\PopnGroFac$ is the population's growth factor and $\lnat$ is the natural log, so $\PopnGroRte = \log \PopnGroFac$.}
\begin{equation}\begin{gathered}\begin{aligned}
  \dot{\LLvl}/\LLvl & =  \PopnGroRte ,
\end{aligned}\end{gathered}\end{equation}
and $\PtyLab$ is an index of labor productivity that grows at rate\marginFootnote{This is ``labor-augmenting'' (Harrod-neutral) productivity growth: $\PtyLab$ multiplies labor, $\FFunc(\KLvl,\PtyLab \LLvl)$.
The alternative is productivity $\PtyAgg$ that multiplies total output, $\PtyAgg \FFunc(\KLvl,\LLvl)$ (``Hicks-neutral'').
With any production function other than Cobb-Douglas, a steady state with constant growth rates requires progress to be labor-augmenting \citep{uzawaNeutral}.}
\begin{equation}\begin{gathered}\begin{aligned}
  \dot{\PtyLab }/\PtyLab  & =  \PtyGroRte.
\end{aligned}\end{gathered}\end{equation}

Technological progress thus lets each worker produce ever more with the same physical capital.

The quantity $\PtyLab \LLvl$ is the number of ``efficiency units'' of labor.

Aggregate capital accumulates according to
\begin{equation}\begin{gathered}\begin{aligned}
  \dot{\KLvl} & =  \YLvl - \CLvl - \deprRte \KLvl \label{eq:Kdotlevel}.
\end{aligned}\end{gathered}\end{equation}

Lower case variables are the upper case version divided by efficiency units,\marginFootnote{NOT per capita; this makes the ``\hyperref[sec:RCK-SteadyState]{balanced growth}'' path a steady state.} so $\yNrm = \YLvl/(\PtyLab \LLvl) = (\KLvl/(\PtyLab \LLvl))^{\kapShare} = \kNrm^{\kapShare}$. Since $\kNrm = \KLvl/(\PtyLab \LLvl)$,
\begin{equation}
  \dot{\kNrm}/\kNrm = \dot{\KLvl}/\KLvl - \dot{\PtyLab}/\PtyLab - \dot{\LLvl}/\LLvl = \dot{\KLvl}/\KLvl - (\PtyGroRte +\PopnGroRte),
\end{equation}
so dividing \eqref{eq:Kdotlevel} by $\PtyLab \LLvl$ and using it gives
\begin{equation}
  \dot{\kNrm} = \fFunc(\kNrm) - \cNrm - (\PtyGroRte +\PopnGroRte+\deprRte) \kNrm . \label{eq:RCK-kdot}
\end{equation}

A steady state is a point where $\dot{\kNrm} = 0$, so that
\begin{equation}
  \cNrm = \fFunc(\kNrm) - (\PtyGroRte +\PopnGroRte+\deprRte) \kNrm . \label{eq:RCK-cSteadyState}
\end{equation}
A first candidate for the best steady state is the one with the highest $\cNrm$. Maximizing steady-state consumption with respect to $\kNrm$ requires
\begin{equation*}
  \frac{d\cNrm}{d\kNrm} = \fFunc^{\prime}(\kNrm) - (\PtyGroRte +\PopnGroRte+\deprRte) = 0 ,
\end{equation*}
and since $d^{2}\cNrm/d\kNrm^{2} = \fFunc^{\prime\prime}(\kNrm) < 0$, this is a maximum. The capital stock that satisfies it, $\GoldRule{\kNrm}$, obeys the ``golden rule'' of \citet{phelps:golden}:\marginFootnote{Phelps called it the ``golden rule'' on the grounds that every generation does unto its successors what it would want its ancestors to have done unto it. You might worry that the argument doesn't make sense, because what is really being equalized across generations is consumption per efficiency unit of labor, and no extant texts of the Christian New Testament quote Jesus as having said ``do unto other efficiency units of labor as you, an efficiency unit of labor yourself, would have liked the ancestral efficiency units to have done unto you.'' Keep reading.}
\begin{equation}
  \fFunc^{\prime}(\GoldRule{\kNrm}) = \PtyGroRte +\PopnGroRte+\deprRte . \label{eq:RCK-goldenRule}
\end{equation}
One more unit of $\kNrm$ adds $\fFunc^{\prime}(\kNrm)$ to output, but holding it steady as effective labor grows and capital depreciates takes $\PtyGroRte +\PopnGroRte+\deprRte$ more investment; consumption is highest where the two are equal. With the Cobb-Douglas $\fFunc(\kNrm) = \kNrm^{\kapShare}$, the golden rule is
\begin{equation}
  \GoldRule{\kNrm} = \left(\kapShare/(\PtyGroRte +\PopnGroRte+\deprRte)\right)^{1/(1-\kapShare)} . \label{eq:RCK-goldenRuleCobbDouglas}
\end{equation}

\section{The Social Planner's Problem}

Now suppose a social planner maximizes the discounted
sum of CRRA utility from per capita consumption, with time preference rate $\discRte$:\marginFootnote{This objective counts each date's per capita utility once, as \citet{blanchard&fischer:text} do. \citet{romer:text} weights it by the size of the population, which replaces $\discRte$ with $\discRte-\PopnGroRte$ below, and takes $\PopnGroRte$ out of the steady-state condition.}
\begin{equation}
  \max \int_{0}^{\infty} \left(\frac{(\CLvl/\LLvl)^{1-\CRRA}}{1-\CRRA}\right) e^{-\discRte  t}\,dt . \label{eq:RCK-maxprob}
\end{equation}

But $\CLvl/\LLvl = \PtyLab (\CLvl/(\PtyLab \LLvl))= \PtyLab \cNrm$.  Recall that for a variable growing at rate $\PtyGroRte$,
\begin{equation}\begin{gathered}\begin{aligned}
  \PtyLab _{t} & =  \PtyLab _{0}e^{\PtyGroRte t} ,
\end{aligned}\end{gathered}\end{equation}
so if productivity at date 0 is $\PtyLab_{0}$, at date $t$
\begin{equation}\begin{gathered}\begin{aligned}
  \CLvl/\LLvl & =  \cNrm\PtyLab   \\
               & =  \cNrm\PtyLab _{0}e^{\PtyGroRte t}. \label{eq:cRat}
\end{aligned}\end{gathered}\end{equation}
Since $\PtyLab_{0}$ multiplies consumption at every date, it only scales the objective, by the positive constant $\PtyLab_{0}^{1-\CRRA}$, so we can impose $\PtyLab_{0}=1$ without loss of generality.

Using \eqref{eq:cRat}, and writing $\uFunc(\cNrm) = \cNrm^{1-\CRRA}/(1-\CRRA)$, with $\CRRA$ the coefficient of relative risk aversion, rewrite the social planner's objective function as
\begin{equation}\begin{gathered}\begin{aligned}
  \int_{0}^{\infty} \left(\frac{(\PtyLab \cNrm)^{1-\CRRA}}{1-\CRRA}\right) e^{-\discRte  t}\,dt & =
  \int_{0}^{\infty} \left(\frac{\cNrm^{1-\CRRA}}{1-\CRRA}\right) e^{-\discRte  t}e^{(1-\CRRA)\PtyGroRte  t}\,dt
  \\ & =  \int_{0}^{\infty} \uFunc(\cNrm) e^{\left((1-\CRRA)\PtyGroRte -\discRte \right) t}\,dt
           .
           \label{eq:maxprobmod}
\end{aligned}\end{gathered}\end{equation}

Thus, assuming $\discRte > (1-\CRRA) \PtyGroRte$, so that the objective is finite (that is, imposing the \href{https://llorracc.github.io/BufferStockTheory/#pffvac}{finite value of autarky condition}), the complete optimization problem
is
\begin{equation}
  \max \int_{0}^{\infty} \uFunc(\cNrm) e^{-(\discRte - (1-\CRRA)\PtyGroRte) t}\,dt
\end{equation}
subject to
\begin{equation}
  \dot{\kNrm} = \kNrm^{\kapShare}-\cNrm-(\PopnGroRte+\PtyGroRte+\deprRte)\kNrm,
\end{equation}
whose current-value Hamiltonian is\marginFootnote{If the Hamiltonian is new to you, Peter Thompson's \href{https://www.andrew.cmu.edu/course/88-737/optimal_control/optimal_control.pdf}{lecture notes on optimal control} derive the method from the economics of a firm's capital problem, and later work through this model. Wikipedia's \href{https://en.wikipedia.org/wiki/Hamiltonian_(control_theory)\#Example:_Ramsey\%E2\%80\%93Cass\%E2\%80\%93Koopmans_model}{worked example of the RCK model} uses the present-value Hamiltonian instead, whose costate is this one discounted to date zero.}
\begin{equation}
  \Ham(\kNrm,\cNrm,\vk) = \uFunc(\cNrm) + (\kNrm^{\kapShare} - \cNrm - (\PopnGroRte+\PtyGroRte+\deprRte)\kNrm)\vk.
\end{equation}

The \href{https://intertemporal-choice.github.io/content/appendices/mathematical/hamiltonianvsdiscrete/\#eq-uceqvk}{current-value Hamiltonian optimization condition} for consumption requires $\partial
\Ham/\partial \cNrm = 0$:
\begin{equation}\begin{gathered}\begin{aligned}
  \cNrm^{-\CRRA} & =  \vk  \label{eq:RCK-H1A} \\
  -\CRRA \cNrm^{-\CRRA-1}\dot{\cNrm} & =  \dot{\vk}
                                       .
\end{aligned}\end{gathered}\end{equation}

The \href{https://intertemporal-choice.github.io/content/appendices/mathematical/hamiltonianvsdiscrete/\#eq-lamgrow}{current-value Hamiltonian optimization condition} for capital, the costate equation, requires:\marginFootnote{In the book's mathematical appendices, \href{https://intertemporal-choice.github.io/content/appendices/mathematical/hamiltonianvsdiscrete/}{Ramsey Growth in Discrete and Continuous Time} derives the Hamiltonian's conditions for a version of this model without productivity growth and shows how each relates to the discrete-time Bellman approach.}
\begin{equation}\begin{gathered}\begin{aligned}
  \dot{\vk} & =  (\discRte-(1-\CRRA)\PtyGroRte)\vk - \partial \Ham/\partial \kNrm  \\
                & =  (\discRte-(1-\CRRA)\PtyGroRte)\vk - \vk (\fFunc^{\prime}(\kNrm) - (\PopnGroRte+\PtyGroRte+\deprRte))  \\
  \dot{\vk}/{\vk} & =  \discRte+\PopnGroRte+\deprRte+\CRRA \PtyGroRte-\fFunc^{\prime}(\kNrm) .
\end{aligned}\end{gathered}\end{equation}

Since \eqref{eq:RCK-H1A} gives $\dot{\vk}/\vk = -\CRRA\dot{\cNrm}/\cNrm$, this is
\begin{equation}
  \dot{\cNrm}/\cNrm = \CRRA^{-1}(\fFunc^{\prime}(\kNrm)-(\PopnGroRte+\deprRte)-\discRte - \CRRA \PtyGroRte), \label{eq:cdotOc}
\end{equation}
or, with $\rfree \equiv \fFunc^{\prime}(\kNrm) - \deprRte$ the interest rate, $\dot{\cNrm}/\cNrm = \CRRA^{-1}(\rfree - \discRte - \PopnGroRte) - \PtyGroRte$.

This Euler equation is called the ``Keynes-Ramsey rule,'' because Ramsey derived it
with an explanation attributed to Keynes.
Apart from the $-\PtyGroRte$, which comes from measuring consumption per efficiency unit, it is like \href{https://intertemporal-choice.github.io/content/consumption/perfforesightcrra/#eq-pfcrra-cgrow}{the Euler equation in the perfect foresight CRRA model} in continuous time, with two differences: the return that matters is $\rfree-\PopnGroRte$ rather than $\rfree$, because population growth dilutes each person's capital, and that return varies over time as $\fFunc^{\prime}(\kNrm)$ does.

Since per capita consumption is $\CLvl/\LLvl = \cNrm \PtyLab$, it grows at $\dot{\cNrm}/\cNrm + \PtyGroRte = \CRRA^{-1}(\fFunc^{\prime}(\kNrm)-(\PopnGroRte+\deprRte)-\discRte)$: the same formula, as a function of $\kNrm$, as in the model with no productivity growth, the \hyperref[eq:cdotOc]{Euler equation} with $\PtyGroRte=0$. Productivity growth matters only through $\kNrm$.

\section{The Steady State}\label{sec:RCK-SteadyState}

Technological progress is assumed to be labor-augmenting because then,
in the steady state, per capita consumption, income
and capital all grow at rate $\PtyGroRte$.\marginFootnote{\citet{ghosDespiteUzawa} discuss how realistic this requirement is.}

$\dot{\cNrm}/\cNrm = 0$ implies that at the steady-state value of $\Steady{\kNrm}$,
\begin{equation}
  \fFunc^{\prime}(\Steady{\kNrm}) = \discRte +\PopnGroRte+\deprRte+\CRRA \PtyGroRte . \label{eq:RCK-modGoldenRule}
\end{equation}

This is the ``modified golden rule'' of \citet{blanchard&fischer:text}, chapter 2. Comparing it with the \hyperref[eq:RCK-goldenRule]{golden rule}, $\fFunc^{\prime}(\Steady{\kNrm}) - \fFunc^{\prime}(\GoldRule{\kNrm}) = \discRte - (1-\CRRA)\PtyGroRte$, so $\Steady{\kNrm} < \GoldRule{\kNrm}$ exactly when $\discRte > (1-\CRRA)\PtyGroRte$: the assumption that keeps the planner's objective finite also rules out accumulating more capital than the golden rule.

Using $\fFunc^{\prime}(\kNrm) = \kapShare \kNrm^{\kapShare-1}$, the modified golden rule gives an analytical formula for the steady state:
\begin{equation}\begin{gathered}\begin{aligned}
  \kapShare \Steady{\kNrm}^{\kapShare-1} & =  \discRte +\PopnGroRte+\deprRte+\CRRA \PtyGroRte  \\
  \Steady{\kNrm} & =  \left(\frac{\kapShare}{\discRte +\PopnGroRte+\deprRte+\CRRA \PtyGroRte}\right)^{\frac{1}{1-\kapShare}} . \label{eq:RCK-kSteadyCobbDouglas}
\end{aligned}\end{gathered}\end{equation}

Thus $\Steady{\kNrm}$ is higher when capital's share of output, $\kapShare$, is higher, and lower when consumers are more impatient, population grows faster, capital depreciates faster, or productivity grows faster.\marginFootnote{Strictly, $\Steady{\kNrm}$ rises with $\kapShare$ when $\kapShare/(\discRte+\PopnGroRte+\deprRte+\CRRA\PtyGroRte) > e^{-(1-\kapShare)/\kapShare}$, which holds for any plausible parameters: with the notebook's, the left side is 1.76 and the right 0.097.} Faster productivity growth lowers it for the reason behind the \href{https://intertemporal-choice.github.io/content/consumption/perfforesightcrra/\#eq-patwgrocond}{growth impatience condition}: consumers who know they will be richer in the future want to spend more today. When productivity grows, it is also lower when $\CRRA$ is higher: consumers who dislike differences in consumption over time save less when they expect their consumption to grow.

\section{A Phase Diagram}

The RCK model's steady state has an analytical solution, but its transition
does not. The usual tool for analyzing it is a \href{https://en.wikipedia.org/wiki/Phase_portrait}{phase diagram} in $\cNrm$ and $\kNrm$.

The first step is to find where the system's differential equations are
zero. By \eqref{eq:RCK-kdot}, $\dot{\kNrm}=0$ implies
\begin{equation}\begin{gathered}\begin{aligned}
  \cNrm & =  \fFunc(\kNrm) - (\PtyGroRte +\PopnGroRte+\deprRte) \kNrm , \label{eq:RCK-kdotEq0}
\end{aligned}\end{gathered}\end{equation}
and the $\dot{\cNrm}/\cNrm=0$ locus is the constant $\Steady{\kNrm}$ found above.
The two loci divide the diagram into its phases, as figure~\ref{fig:RamseySSPlot} shows.

\begin{figure}
  \caption{$\dot{\cNrm}/\cNrm =0 $ and $\dot{\kNrm}=0$ Loci}\label{fig:RamseySSPlot}
  \includegraphics[width=6in]{\FigDir/RamseySSPlot.png}

  \small Note: With $\CRRA=2$, $\kapShare=0.3$, $\discRte=0.02$, $\PopnGroRte=0.01$ and $\deprRte=0.08$, and productivity growth $\PtyGroRte$ of 0.03 (low) or 0.06 (high); the dotted line marks the golden rule when $\PtyGroRte$ is low. Drawn by the \href{\Notebook}{notebook}; run it in your browser with nothing to install, in \href{\BookNotebookLite}{JupyterLite}, or open it in \href{\NotebookColab}{Colab} or \href{\NotebookBinder}{Binder}.
\end{figure}

\section{Transition}

The two differential equations and $\kNrm_{0}$ don't pin down $\cNrm_{0}$: every starting $\cNrm_{0}$ traces a path. Together with the requirement that capital stay positive, the transversality condition picks one:\marginFootnote{The transversality condition is an optimality condition, not the intertemporal budget constraint. The requirement that capital stay positive, $\kNrm > 0$, rules out ending with negative wealth in present value; the transversality condition rules out ending with positive wealth in present value, capital that is never consumed. Together they make the intertemporal budget constraint hold with equality. \href{https://intertemporal-choice.github.io/content/appendices/mathematical/hamiltonianvsdiscrete/}{Ramsey Growth in Discrete and Continuous Time} shows that the condition sets the present value of capital in the far future to zero, in discrete and in continuous time.}
\begin{equation}\begin{gathered}\begin{aligned}
  \lim_{t \rightarrow \infty} e^{-\int_{0}^{t} (\rfree_{\tau} -\PopnGroRte -\PtyGroRte)\,d\tau }\kNrm_{t} & =  0. \label{eq:tvc}
\end{aligned}\end{gathered}\end{equation}

Figure~\ref{fig:RamseySaddlePlot} shows three paths for $\cNrm$ and $\kNrm$ that satisfy the dynamic equations for consumption and assets, \eqref{eq:cdotOc} and \eqref{eq:RCK-kdot}, all from the same $\kNrm_{0}$. The path that starts above the ``saddle path'' (the middle one) runs capital to zero in finite time, after which consumption would have to collapse, which the Euler equation forbids. The path that starts below it accumulates capital past the golden rule, toward the point where the $\dot{\kNrm}=0$ locus meets the $\kNrm$ axis, with consumption falling toward zero; there the interest rate is below $\PopnGroRte+\PtyGroRte$, so the discount term grows without bound and the \hyperref[eq:tvc]{transversality condition} fails. Only the saddle path satisfies all three conditions.\marginFootnote{The steady state is a saddle point. Linearize the system there: its Jacobian matrix, the derivatives of $\dot{\kNrm}$ and $\dot{\cNrm}$ with respect to $\kNrm$ and $\cNrm$, has $\partial\dot{\kNrm}/\partial\cNrm = -1$ and $\partial\dot{\cNrm}/\partial\cNrm = 0$, so its determinant is $\partial\dot{\cNrm}/\partial\kNrm = \Steady{\cNrm}\fFunc^{\prime\prime}(\Steady{\kNrm})/\CRRA < 0$, and its two eigenvalues have opposite signs: one direction converges to the steady state, and the other diverges. The saddle path is the converging one; the stable eigenvalue, in absolute value, is the model's speed of convergence, which \citet{barro&salaimartin:book}, chapter 2, compare with the speed observed across countries and regions.}

\begin{figure}
  \caption{Transition to the Steady State}\label{fig:RamseySaddlePlot}
  \includegraphics[width=6in]{\FigDir/RamseySaddlePlot.png}

  \small Note: The same economy, with $\PtyGroRte=0.03$. From $\kNrm=0.5$, consumption 10 percent above the saddle path runs the capital stock down to zero, and consumption 10 percent below it accumulates capital past the golden rule, marked by the dotted line. Drawn by the \href{\Notebook}{notebook}; run it in your browser with nothing to install, in \href{\BookNotebookLite}{JupyterLite}, or open it in \href{\NotebookColab}{Colab} or \href{\NotebookBinder}{Binder}.
\end{figure}

The saddle path has, in general, no analytical solution, so the transition must be computed. The \href{\Notebook}{notebook} computes it by the ``time elimination'' method \citep{Mulligan&sm:transdynam}: along the path, $d\cNrm/d\kNrm = \dot{\cNrm}/\dot{\kNrm}$, a differential equation in $\kNrm$ alone, which it solves outward from the steady state, starting along the saddle path's slope there. The notebook derives the method and draws both figures; run it in your browser with nothing to install, in \href{\BookNotebookLite}{JupyterLite}, or open it in \href{\NotebookColab}{Colab} or \href{\NotebookBinder}{Binder}. \citet{judd:book} covers numerical methods for saddle paths more generally.

How much of observed growth can the transition explain? Little, \citet{king&rebelo:trans} argue: with plausible parameters, an economy far enough below its steady state for the transition to account for much of its growth would have implausibly high interest rates early on.

\href{https://intertemporal-choice.github.io/content/growth/decentralizingrck/}{Decentralizing the Ramsey/Cass-Koopmans Model} shows that, under special conditions, competitive households and firms, taking prices as given, reproduce the planner's allocation, and \href{https://intertemporal-choice.github.io/content/growth/rckwithgov/}{The Ramsey/Cass-Koopmans (RCK) Model with Government} adds government spending, taxes and debt.

\bibliographystyle{econark}
\bibliography{\bibfilesfound}

\end{document}
